% Chapter 3: Problem Definition & Requirements Engineering
% BSU Engineering Portal Graduation Project

\chapter{Problem Definition \& Requirements Engineering}
\label{chap:requirements}

\section{Current Situation \& Stakeholder Analysis}
\label{sec:current-situation}

\subsection{The Pre-System Workflow}

Prior to the introduction of the BSU Engineering Portal, the Faculty of Engineering at Beni-Suef University operated a predominantly paper-based and spreadsheet-driven academic management process. This workflow exhibited the following systemic failures:

\begin{itemize}
    \item \textbf{Student Records}: Student enrollment data was maintained in Microsoft Excel workbooks, segregated by department and academic year. Adding a single student required manually editing a shared file over the local network, risking data corruption due to concurrent edits.
    \item \textbf{Grade Management}: Faculty members (doctors) recorded attendance and exam grades on paper forms. These forms were physically transported to Student Affairs, where clerks manually re-entered them into another Excel workbook---introducing a high-latency human transcription layer with zero automated validation.
    \item \textbf{Course-Doctor Assignment}: The assignment of faculty members to courses each semester was communicated verbally or via internal paper memos, lacking a central digital record mapping doctors to subjects in the current term.
    \item \textbf{Student Access to Information}: Students lacked a self-service mechanism to view their own grades, attendance records, or exam schedules. They depended entirely on physical notice boards or direct inquiries to Student Affairs, processes that could take days.
    \item \textbf{Audit Trail}: The absence of digital logs meant there was no record of who uploaded grades, who modified student records, or when these actions occurred, making accountability impossible and dispute resolution arbitrary.
\end{itemize}

\subsection{Stakeholder Identification}

Identifying the exact users of the system is the first step in requirements engineering. Table~\ref{tab:stakeholders} identifies every actor in the system, their organizational role, their pain points under the legacy process, and their primary interaction with the proposed portal.

\begin{figure}[H]
    \centering
    \includegraphics[width=\textwidth]{PLACEHOLDER_STAKEHOLDER_DIAGRAM.png}
    \caption{PLACEHOLDER: Insert a Stakeholder Interaction Diagram showing all five roles (Admin, Doctor, Student, Student Affairs, Staff Affairs) and their primary data flows into the centralized system.}
    \label{fig:stakeholders}
\end{figure}

\begin{table}[H]
    \centering
    \renewcommand{\arraystretch}{1.3}
    \begin{tabular}{|p{0.15\textwidth}|p{0.2\textwidth}|p{0.3\textwidth}|p{0.25\textwidth}|}
        \hline
        \textbf{Stakeholder} & \textbf{Organization Role} & \textbf{Pre-System Pain Point} & \textbf{System Interaction} \\
        \hline
        \textbf{Student} & Undergraduate in any dept/level & No self-service access to grades, attendance, or materials. & View grades, attendance, quizzes, schedule, and download lectures. \\
        \hline
        \textbf{Doctor} & Faculty Member teaching courses & No digital tool for recording attendance or publishing quizzes. & Manage courses, record attendance, create quizzes, enter grades. \\
        \hline
        \textbf{Student Affairs} & Admin staff managing student lifecycles & Manual Excel-based registration and redundant grade re-entry. & Batch-upload students, view grades, manage certificates. \\
        \hline
        \textbf{Staff Affairs} & Admin staff managing HR & No structured system for managing doctor course assignments. & Batch-upload doctors, assign to courses, manage deletions. \\
        \hline
        \textbf{Admin} & IT/System Administrator & No central configuration tool; ad-hoc management. & Manage academic years, approve deletions, view audit logs. \\
        \hline
    \end{tabular}
    \caption{Stakeholder Summary and System Interaction Points}
    \label{tab:stakeholders}
\end{table}

\subsection{Stakeholder Permission Boundaries}

A critical engineering decision in the system's design is the enforcement of the \textbf{principle of least privilege}. No single stakeholder role has full system access. Even the Admin role cannot directly enter grades or upload students---those operational capabilities are delegated exclusively to the Doctor and Student Affairs roles, respectively. This ensures that operational authority perfectly matches organizational responsibility.

The following code snippet demonstrates how these boundaries are strictly enforced at the API layer using Django REST Framework permission classes:

\begin{lstlisting}[language=Python, caption=Role-based Permission Guard for Student Affairs (\texttt{permissions.py})]
class IsStudentAffairsRole(permissions.BasePermission):
    """Student Affairs role - exclusively manages students and grades"""
    def has_permission(self, request, view):
        if not request.user or not request.user.is_authenticated:
            return False
        # Admin gets fallback read/write access for administrative support
        return request.user.role in ['STUDENT_AFFAIRS', 'ADMIN']
\end{lstlisting}

\section{Proposed Methodology}
\label{sec:methodology}

\subsection{Why Agile?}

The BSU Engineering Portal was developed using an \textbf{iterative, sprint-based Agile methodology}. A traditional Waterfall approach would have mandated locking all requirements prior to implementation. However, stakeholder interviews frequently revealed new requirements throughout the development cycle. For example, the necessity for a multi-modal quiz engine (supporting MCQ, Essay, and Image-based questions) emerged mid-development when faculty expressed frustration over the lack of integrated online assessment tools.

Agile permitted the engineering team to:
\begin{enumerate}
    \item \textbf{Deliver working features early}: Students could view grades while the quiz engine was still being actively developed.
    \item \textbf{Incorporate feedback incrementally}: The grading template logic was fundamentally redesigned after pilot testing revealed the initial model did not accommodate the faculty's official tertiary grading structure (Coursework / Written / Practical-Oral).
    \item \textbf{Maintain continuous deployment}: Docker containerization ensured every sprint concluded with a highly reproducible, runnable artifact.
\end{enumerate}

\begin{table}[H]
    \centering
    \renewcommand{\arraystretch}{1.2}
    \begin{tabular}{|p{0.15\textwidth}|p{0.3\textwidth}|p{0.45\textwidth}|}
        \hline
        \textbf{Sprint} & \textbf{Focus Area} & \textbf{Key Deliverables} \\
        \hline
        Sprint 1 & Auth \& Core Data & 5-Role User model, login/logout, CSRF, Core Academic Models. \\
        Sprint 2 & Student \& Staff Affairs & Batch Excel uploads (Pandas), student filtering. \\
        Sprint 3 & Course Management & CourseOffering model, doctor assignments, lecture uploads. \\
        Sprint 4 & Attendance \& Grading & Bulk attendance recording, grading templates. \\
        Sprint 5 & Quiz Engine & MCQ/Essay CRUD, auto-grading logic. \\
        Sprint 6 & Student Portal & Dashboards, schedules, grade views. \\
        Sprint 7 & Admin \& Operations & Announcements, audit logging, deletion workflows. \\
        Sprint 8 & DevOps \& Polish & Docker Compose, Nginx proxy, RTL (Arabic) UI support. \\
        \hline
    \end{tabular}
    \caption{Agile Sprint Plan Overview}
    \label{tab:sprints}
\end{table}

\section{Functional Requirements}
\label{sec:functional-reqs}

The following functional requirements are classified using IEEE 830-style prioritization (Must / Should / Could). 

\subsection{Authentication \& User Management}

\begin{itemize}
    \item \textbf{FR-01 [Must]}: The system shall authenticate users via session-based login using a username (national ID) and password.
    \item \textbf{FR-02 [Must]}: The system shall enforce a mandatory password change on first login. The default password shall equal the user's national ID.
    \item \textbf{FR-03 [Must]}: The system shall reject new passwords that match the national ID or are under 6 characters.
    \item \textbf{FR-04 [Must]}: The system shall enforce isolated UI dashboards for all five defined user roles.
    \item \textbf{FR-05 [Must]}: Admin and authorized personnel shall be able to reset any user's password, automatically re-triggering the first-login requirement flag.
\end{itemize}

\begin{lstlisting}[language=Python, caption=Enforcing First-Login Requirement in Model Save (\texttt{models.py})]
def save(self, *args, **kwargs):
    is_new = not self.pk
    if is_new and not self.is_superuser:
        if not self.password or self.password == '!':
            self.set_password(self.national_id or self.username)
        if self.role in ['STUDENT', 'STUDENT_AFFAIRS', 'STAFF_AFFAIRS', 'DOCTOR']:
            self.first_login_required = True
    return super().save(*args, **kwargs)
\end{lstlisting}

\subsection{Student Affairs Operations}

\begin{itemize}
    \item \textbf{FR-06 [Must]}: Upload students in bulk via Excel/CSV (national\_id, full\_name, department, level).
    \item \textbf{FR-07 [Should]}: Provide a pre-flight UI preview of the first 5 rows containing validation errors before committing database writes.
    \item \textbf{FR-08 [Must]}: Validate that Excel data matches dropdown selections (e.g., rejecting a "Mechanical" student uploaded into the "Electrical" form).
    \item \textbf{FR-09 [Must]}: Perform an "upsert" during batch uploads---updating existing records rather than crashing on duplicate national IDs.
    \item \textbf{FR-10 [Must]}: View all student grades globally in a read-only matrix grouped by student and subject.
    \item \textbf{FR-11 [Must]}: Upload zipped PDF graduation certificates, auto-matching filenames (e.g., \texttt{29912...pdf}) to student records.
    \item \textbf{FR-12 [Must]}: Maintain a permanent \texttt{UploadHistory} summarizing created/updated/failed rows per batch.
\end{itemize}

\begin{figure}[H]
    \centering
    \includegraphics[width=\textwidth]{PLACEHOLDER_UPLOAD_PREVIEW.png}
    \caption{PLACEHOLDER: Insert a screenshot of the Student Affairs Upload Form showing the 5-row preview validation step highlighting row errors.}
    \label{fig:upload_preview}
\end{figure}

\subsection{Staff Affairs \& Doctor Operations}

\begin{itemize}
    \item \textbf{FR-13 [Must]}: Staff Affairs shall batch-upload Doctors and assign them to specific \texttt{CourseOfferings} per term.
    \item \textbf{FR-14 [Must]}: Staff Affairs shall submit Doctor deletion requests to the Admin rather than deleting them directly, preventing accidental cascading data loss.
    \item \textbf{FR-15 [Must]}: Doctors shall record bulk attendance for course sessions and export reports.
    \item \textbf{FR-16 [Must]}: Doctors shall enter bulk grades mapped strictly to the course's \texttt{GradingTemplate}.
    \item \textbf{FR-17 [Must]}: Doctors shall create Quizzes (MCQ, Essay, Image). MCQ attempts shall be auto-graded; Essay attempts shall queue for manual doctor review.
\end{itemize}

\section{Non-Functional Requirements}
\label{sec:non-functional-reqs}

Non-functional requirements dictate the system's operational constraints and quality attributes.

\begin{table}[H]
    \centering
    \renewcommand{\arraystretch}{1.3}
    \begin{tabular}{|p{0.1\textwidth}|p{0.15\textwidth}|p{0.45\textwidth}|p{0.2\textwidth}|}
        \hline
        \textbf{ID} & \textbf{Category} & \textbf{Requirement} & \textbf{Target Metric} \\
        \hline
        \textbf{NFR-01} & Security & Passwords hashed via PBKDF2-SHA256. & 0 plaintext DB records \\
        \textbf{NFR-02} & Security & API enforces CSRF protection with \texttt{SameSite=None} cookies for cross-origin integration. & 100\% state mutations verified \\
        \textbf{NFR-03} & Performance & API responses for list endpoints complete in under 2s for datasets $\leq$ 5,000 rows. & $<2$s response time \\
        \textbf{NFR-04} & Usability & Full RTL layout support for Arabic text via MUI plugins. & No broken RTL layouts \\
        \textbf{NFR-05} & Audit & All administrative mutations logged in \texttt{AuditLog}. & 100\% audit coverage \\
        \hline
    \end{tabular}
    \caption{Non-Functional Requirements}
    \label{tab:nfrs}
\end{table}

To protect against brute-force and denial-of-service attempts, rate limiting was configured globally in the backend:

\begin{lstlisting}[language=Python, caption=Global API Rate Limiting Configuration (\texttt{settings.py})]
REST_FRAMEWORK = {
    'DEFAULT_THROTTLE_CLASSES': [
        'rest_framework.throttling.AnonRateThrottle',
        'rest_framework.throttling.UserRateThrottle',
    ],
    'DEFAULT_THROTTLE_RATES': {
        'anon': '100/hour',    # Protects login/public endpoints
        'user': '1000/hour',   # Accommodates heavy batch operations
    },
}
\end{lstlisting}

\section{Use Case Definitions}
\label{sec:use-cases}

\begin{figure}[H]
    \centering
    \includegraphics[width=\textwidth]{PLACEHOLDER_USE_CASE_UML.png}
    \caption{PLACEHOLDER: Insert a UML Use Case Diagram mapping the 5 actors to their respective grouped use cases (Authentication, Academic Mgmt, Quizzes, etc.).}
    \label{fig:use_case_uml}
\end{figure}

\subsection{Detailed Use Case: Batch Student Upload}

This use case outlines the system's most complex data validation pipeline, migrating the faculty away from disjointed Excel hand-offs. A critical engineering decision was made to process rows iteratively inside \texttt{transaction.atomic()} blocks. While Django's \texttt{bulk\_create()} is faster, it fails entirely if a single row errors and obscures row-level validation feedback. The iterative atomic approach allows valid rows to succeed while cleanly capturing exact row indices for failed records.

\textbf{Primary Actor:} Student Affairs\\
\textbf{Main Flow:}
\begin{enumerate}
    \item User selects Department, Academic Year, and Level, then uploads a CSV/Excel file.
    \item System parses the file via Pandas, validating headers and verifying that the file's payload matches the UI dropdown selections (preventing cross-department contamination).
    \item System iterates over rows. If a \texttt{national\_id} exists, the record is updated. If not, a new \texttt{User} and \texttt{Student} record are atomically created.
    \item The system commits an \texttt{UploadHistory} audit log and returns a JSON summary.
\end{enumerate}

\begin{figure}[H]
    \centering
    \includegraphics[width=\textwidth]{PLACEHOLDER_SEQUENCE_UPLOAD.png}
    \caption{PLACEHOLDER: Insert a Sequence Diagram showing the Batch Student Upload flow (Frontend $\rightarrow$ API $\rightarrow$ Pandas Validation $\rightarrow$ Database Transaction $\rightarrow$ AuditLog).}
    \label{fig:seq_upload}
\end{figure}

\subsection{Detailed Use Case: Doctor Deletion Approval Workflow}

Deleting a doctor account cascades destructively through the relational database, removing \texttt{CourseOfferings}, \texttt{StudentGrades}, and \texttt{Attendance} records. To mitigate the risk of catastrophic data loss by administrative error, a \textbf{two-phase approval pattern} was implemented.

\textbf{Main Flow:}
\begin{enumerate}
    \item Staff Affairs initiates a deletion request, providing a required justification.
    \item The API generates a \texttt{DoctorDeletionRequest} marked as \texttt{PENDING}. No data is deleted.
    \item The Admin reviews the request on the Admin Dashboard.
    \item If approved, the Admin triggers the hard deletion and an \texttt{AuditLog} is captured. If rejected, the request is closed and the account remains untouched.
\end{enumerate}

\begin{figure}[H]
    \centering
    \includegraphics[width=0.8\textwidth]{PLACEHOLDER_ADMIN_APPROVAL.png}
    \caption{PLACEHOLDER: Insert a screenshot of the Admin Approval Center showing pending doctor deletion requests with Approve/Reject actions.}
    \label{fig:admin_approval}
\end{figure}

\subsection{Detailed Use Case: Quiz Lifecycle}

The quiz engine represents a fully integrated state machine tracking an assessment from creation to final grading.

\textbf{Main Flow:}
\begin{enumerate}
    \item Doctor provisions a quiz, specifying MCQ or Essay questions and a time limit.
    \item Student starts an attempt; system marks status as \texttt{IN\_PROGRESS}.
    \item Student submits. For MCQ questions, the system instantly calculates scores by comparing choices against the \texttt{is\_correct} boolean flag.
    \item For Essay questions, the attempt is flagged as \texttt{SUBMITTED}. The Doctor must manually review the text payload and assign \texttt{points\_earned}, pushing the status to \texttt{GRADED}.
\end{enumerate}

\begin{lstlisting}[language=Python, caption=Real-time Auto-grading Logic for MCQ Quizzes (\texttt{quiz\_views.py})]
if question.question_type == 'MCQ':
    choice = QuizChoice.objects.get(id=choice_id, question=question)
    answer.selected_choice = choice
    if choice.is_correct:
        answer.points_earned = question.points  # Award full points
        total_score += float(question.points)
    else:
        answer.points_earned = 0                # Incorrect choice
\end{lstlisting}